% Korean Noether paper 39 — fresh producer translation, unchecked.
% Initial translation authority: ED0005 lines 19052--19141.
% Current scope-rebase authority: ED0011 / NOETH-DE-AUTH-v047-20260808.
% Authority SHA-256: 1A44F967B29972E8F99E5C323A479162AD82A23FC457395915A4BB9DDF51AD41.
% Normalized source span: 18,724 bytes; SHA-256 4F6355189925F249DE27FE5FD25C22FB3A2226088EBB7CAF5CB486607A112B7C.
% Translation only; independent source, Korean-language, render, and certification checks remain external.

\section*{39. 초복소계와 가환대수 및 수론의 관계}

\begin{center}
\emph{Verhandl. Intern. Math.-Kongreß Zürich 1 (1932), 189--194쪽}
\end{center}

\begin{center}
\textsc{괴팅겐의 에미 뇌터}
\end{center}

\begingroup
\renewcommand{\thefootnote}{\arabic{footnote})}
\setcounter{footnote}{0}

\paragraph{1.}
초복소계, 곧 대수의 이론은 최근 몇 년 동안 크게 발전하였다. 그러나 이 이론이 가환적인 문제들에 갖는 의의가 분명해진 것은 극히 최근의 일이다. 오늘 나는 비가환적인 것이 가환적인 것에 갖는 이 의의를 보고하고자 한다. 구체적으로는 가우스까지 거슬러 올라가는 두 고전적 문제, 곧 주종정리와 그에 밀접하게 연결된 노름정리를 따라 이를 살펴보겠다. 이 문제들의 정식화는 시대에 따라 거듭 달라졌다. 가우스에게서는 이차형식론의 결말로 나타났고, 뒤에는 유체론을 통하여 상대 순환 수체와 상대 아벨 수체를 특성화하는 데 본질적인 구실을 하였다. 마지막으로는 자기동형사상과 대수의 분할에 관한 정리로 표현할 수 있으며, 이 마지막 정식화는 동시에 그 정리들을 임의의 상대 갈루아 수체로 옮겨 준다.

뒤에서 자세히 전개할 이 개요를 통하여, 비가환적인 것을 가환적인 것에 적용하는 \emph{원리}도 설명하고자 한다. \emph{대수론을 이용하여 이차형식이나 순환체에 관한 알려진 사실들을 불변적이고 단순한 형태, 곧 대수의 구조적 성질에만 의존하는 형태로 정식화하려 한다. 이러한 불변적 정식화를 일단 증명하면 --- 위에서 든 예들에서는 실제로 증명되어 있다 --- 그 사실들은 저절로 임의의 갈루아체로 옮겨진다.}

\paragraph{2.}
세부를 전개하기 전에 여러 방법과 그 밖의 결과들을 전체적으로 개관하겠다. 우선 가장 큰 어려움은 일반 갈루아체에 맞는 정식화를 얻는 데 있으며, 초복소적 방법이 없었다면 이를 시작할 실마리조차 없었음을 말해 두어야 한다. 앞서 든 예에서 비가환적인 것으로 넘어가는 바탕은 \glqq{}교차곱\grqq{}과 그 곱셈상수인 \glqq{}인자계\grqq{}를 이용하여 \emph{체와 군을 동시에 고찰}하는 것이다(3절 참조). 이렇게 하면 기초체 위의 단순 정규대수 하나를 얻으며, 본질적으로 모든 그러한 대수를 이 방식으로 생성할 수 있다. 이러한 교차곱은 Dickson이 처음 고찰하였다.\footnote{그의 저서 \emph{Algebren und ihre Zahlentheorie}, Zürich 1927, § 34를 보라.} 한편 인자계의 이론은 Speiser, Schur, R. Brauer가 절대기약표현의 문제라는 전혀 다른 출발점에서 발전시켰다.\footnote{예를 들어 R. Brauer, \emph{Untersuchungen über die arithmetischen Eigenschaften von Gruppen linearer Substitutionen}, Math. Z. 28 (1928) 및 그 논문의 주 2)에 제시된 문헌을 보라.} 두 이론을 합친 뒤에야 가환적인 문제를 다룰 만큼 단순하면서도 넓게 적용되는 구성을 얻을 수 있었다.\footnote{나는 이 구성을 1929/30년 겨울학기 강의에서 처음 전개하였다. 그 내용은 H. Hasse, \emph{Theory of cyclic algebras}, Transac. 134 (1932), 제2장에 실려 있다. 이 강연이 다루는 전 분야에 관해서는 M. Deuring의 초복소수와 수론적 응용에 관한 보고가 길잡이가 되며, 이는 \glqq{}Ergebnisse der Mathematik\grqq{} 총서에 출판될 예정이다.}

동시에 알려진 사실들에도 새롭고 투명한 증명이 생긴다. 여기서는 H. Hasse가 제시하여 곧 Math. Ann.에 실릴 순환체의 \emph{상호법칙}에 대한 초복소적 증명을 들고자 한다.\footnote{H. Hasse, \emph{Die Struktur der R. Brauerschen Algebrenklassengruppe über einem algebraischen Zahlkörper (insbesondere Normenrestsymbol und Reziprozitätsgesetz)}, Math. Ann. 107 (1932/33).} 이 증명은 교차곱 이론에 바탕을 둔 그의 노름잉여기호의 불변적 정식화를 이용한다. 또 C. Chevalley가 최근 같은 토대 위에서 제시한 국소 유체론의 초복소적 기초도 있다. 다만 이를 위해서는 인자계에 관한 새로운 대수적 정리들을 더 발전시켜야 했다.\footnote{C. Chevalley, \emph{Sur la théorie du symbole de restes normiques}; J. f. Math. 169에 실릴 예정이다.} 그러나 교차곱의 방법만으로는 아무래도 \emph{갈루아 수체의 이론 전체}를 얻지 못한다는 제한도 덧붙여야 한다. 위에서 말한 Hasse의 증명을 앞의 원리에 따라 잇는 Artin의 아직 출판되지 않은 새 결과들은 완전한 동형정리 대신 개수의 등식만을 주기 때문이다.

완전한 동형, 다만 연산자 동형사상을 주는 방법은 대수적인 범위 안에 이미 존재한다. 이는 A. Speiser의 착상을 이어 간 것이다.\footnote{A. Speiser, \emph{Gruppendeterminante und Körperdiskriminante}, Math. Ann. 77 (1916).} 갈루아체를 \glqq{}갈루아 가군\grqq{}, 곧 기초체 위의 가군으로서 갈루아군의 치환들을 연산자로 허용하는 가군으로 보는 관점이다. 체와 군환(군대수) 사이에는 다음 의미의 연산자 동형사상이 존재한다. 원소들이 일대일로 대응하고 기초체 위의 선형형식들이 서로 대응하며, 체에서의 갈루아군 치환에는 군환에서의 곱셈이 대응한다. 내가 세운 이 정리\footnote{E. Noether, \emph{Normalbasis bei Körpern ohne höhere Verzweigung}, 정리 3, J. f. Math. 167 (1932). 그 증명에는 틈이 있다.}는 M. Deuring이 증명하였다.\footnote{M. Deuring, \emph{Galoissche Theorie und Darstellungstheorie}, Math. Ann. 107 (1932).} 그는 이를 바탕으로 갈루아 이론을 증명하였고, 이때 연산자 동형사상이 군과 체의 대응을 실현한다. 역시 Deuring이 얻은 더 깊은 구조정리들은 Artin \(L\)-급수의 형식적 사실들과 나란히 놓이며 Artin 도체에 구조적으로 접근하는 길을 준다. 일반 군지표로 만드는 이 Artin \(L\)-급수와 도체\footnote{E. Artin, \emph{Über eine neue Art von \(L\)-Reihen}, Math. Sem. Hamburg 3 (1924); \emph{Zur Theorie der \(L\)-Reihen mit allgemeinen Gruppencharakteren}, 같은 곳 8 (1931); \emph{Die gruppentheoretische Struktur der Diskriminanten algebraischer Zahlkörper}, J. f. M. 164 (1931).}는 Speiser의 연구\textsuperscript{6)}와 더불어 수론과 표현론 사이의 첫 연결이며, 아벨체를 넘어선 첫 진전이다. 이 연구들은 전체 발전에 강한 자극을 주었고, 특히 갈루아 가군론은 여기에서 방향을 얻었다.

\paragraph{3.}
이제 처음에 내세운 문제들, 곧 노름정리와 주종정리를 자세히 살펴보겠다. 먼저 \emph{교차곱의 정의}를 제시한다. \(K/k\)를 \(n\)차 갈루아체라 하고 그 군을 \(\Gg\)라 하자. 교차곱이란 \(K\)와 \(\Gg\)를 하나의 대수 \(A\) 안에 동시에 매장하여 \(K\)의 자기동형사상들이 내부자기동형사상이 되게 하는 것이다. \(n\)개의 군 원소에 대응하는 기호들을 \(u_{S_1},\ldots,u_{S_n}\)이라 하자. 그러면 먼저 \(A\)를 \(K\) 위에서 랭크 \(n\)인 선형형식 가군으로 둔다.
\srcnumdisplay[2.8em]{(1)}{\begin{gathered}
A=u_{S_1}K+\cdots+u_{S_n}K\\
\text{\normalfont(즉 \(A\)는 \(a_i\)가 \(K\)의 임의의 원소인}\\
\text{\normalfont 모든 선형형식 \(u_{S_1}a_1+\cdots+u_{S_n}a_n\)으로 이루어진다.)}
\end{gathered}}
내부자기동형사상 --- \(u_S\), 더 일반적으로는 \(u_SK^\ast\)가 생성한다\footnote{\(K^\ast\)는 \(K\)에서 0을 뺀 집합이며, 이하 이 표기를 계속 사용한다.} --- 을 요구하면 \(A\)는 환, 따라서 \(k\) 위에서 랭크 \(n^2\)인 대수가 된다. 그 요구는 모든 \(z\in K\)에 대하여 다음과 같이 표현된다.
\srcnumdisplay[2.8em]{(2)}{u_S^{-1}zu_S=z^S\footnotemark,\qquad \text{또는}\qquad zu_S=u_Sz^S}
\footnotetext{\(z^S\)는 통상적인 표기대로 치환 \(S\)로 \(z\)에서 생기는 원소를 뜻한다.}
\srcnumdisplay[2.8em]{(3)}{u_Su_T=u_{ST}a_{S,T}\quad \text{\(a_{S,T}\in K^\ast\)}}
\srcnumdisplay[2.8em]{(4)}{\begin{gathered}
a_{ST,R}\,a_{S,T}^{\,R}=a_{S,TR}\,a_{T,R}\\
\text{\normalfont(\([u_Su_T]u_R=u_S[u_Tu_R]\)에서 오는 결합법칙)}
\end{gathered}}
\(A\)를 인자계 \(a_{S,T}\)에 대한 \(K\)와 \(\Gg\)의 교차곱이라 한다. \(A\)는 \(k\) 위의 단순 정규대수, 곧 대응하는 나눗셈대수 \(D\) 위의 \(r\)차 행렬환 \(D_r\)이 되고, \(K\)는 극대 가환 부분체, 따라서 분해체가 됨을 증명할 수 있다. 여기서 분해체라는 말은 계수영역 \(k\)를 \(K\)와 동형인 체로 확대하면 대수가 \emph{분할}, 곧 그 중심 위의 완전행렬환이 된다는 뜻이다. 역으로 주어진 나눗셈대수 \(D\)에 대해서도 위와 같은 방식으로 교차곱으로 생성할 수 있는 행렬환 \(D_r\)이 언제나 존재한다.

\(c_S\in K^\ast\)를 택하여 \(u_S\)에서 같은 자기동형사상을 생성하는 \(v_S=u_Sc_S\)로 옮기면 \glqq{}연관된\grqq{} 인자계가 생긴다.
\srcnumdisplay[2.8em]{(5)}{\bar a_{S,T}=a_{S,T}\frac{c_S^{\,T}c_T}{c_{ST}}}
연관된 인자계들은 한 류 \((a)\)로 묶는다. 마찬가지로 \(A\)와 상사인 모든 대수, 곧 \(r=1,2,\ldots\)인 모든 \(D_r\)를 한 류 \(\mathfrak A\)로 묶는다. \emph{류 \(\mathfrak A\)와 \((a)\)는 일대일로 대응한다. 분해체 \(K\)를 고정한 류들은 직접곱 구성에 관하여 아벨군을 이루며, 이는 인자계의 류들을 성분별로 곱하여 얻는 군과 동형이다. 단위원은 각각 분할대수류와 모든 변환량 \(\frac{c_S^{\,T}c_T}{c_{ST}}\)의 계이다.}
이것이 R. Brauer가 발견한 대수류군이다.

\paragraph{4.}
이제 \emph{순환 분해체}로 \emph{특수화}하여 \emph{노름개념}과의 관계, 따라서 처음에 말한 원리에 따른 일반화된 \emph{노름정리}의 정식화에 이르겠다. \(Z\)가 순환체이고 \(S\)가 그 군의 생성치환이면, 대응하는 대수를 순환대수라 한다. \(S\)의 거듭제곱들에 \(u\)의 거듭제곱들을 대응시켜 다음과 같이 둘 수 있다.
\[
\begin{aligned}
(1')\qquad &A=Z+uZ+\cdots+u^{n-1}Z,\\
(2')\qquad &zu=uz^S,\\
(3')\quad &u^n=a,\\
(4')\quad &a\text{는 기초체 }k^\ast\text{에 속한다},\\
(5')\quad &v=uc\text{로 두면 }\bar a=a\cdot N(c).
\end{aligned}
\]
따라서 여기서 각 인자계는 기초체에 속하는 단 하나의 원소 \(a\)로 이루어지며 --- \(A=(a,Z)\)로 표기한다 --- 인자계의 단위원류는 \(Z^\ast\)에서 오는 \emph{노름}들로 주어진다. 대수류군은 몫군 \(k^\ast/N(Z^\ast)\)와 동형이다.
그러므로 순환대수 \((a,Z)\)가 분할되는 것은 \(a\)가 \(Z\)의 한 원소의 노름일 때 그리고 그때뿐이다.

노름과 분할 사이의 이 관계는 \glqq{}노름정리\grqq{}를 다음과 같이 정식화한다. \emph{분할대수에 관한 정리: 한 대수가 모든 소에서 분할되면 전역적으로도 분할된다.} 여기서 \glqq{}소\grqq{}는 수론에서 통상적인 방식으로, 기초체 \(k\)를 그 \(\mathfrak p\)-진 확대 \(k_{\mathfrak p}\)로 바꾸어 정의한다. \(\mathfrak p\)는 \(k\)의 소아이디얼이며, 유한 개의 무한소에 대해서는 \(k\)와 그 켤레체들을 실수체로 확대한다.

실제로 이 명제에는 순환체의 노름정리가 들어 있다. 앞서 본 바에 따라 순환대수 \((a,Z)\)에 대해서는 이 정리가 다음 명제와 동치이기 때문이다. \(a\)가 모든 유한소와 무한소에서 \(\mathfrak p\)-진 노름이면, \(a\)는 \(Z\)의 한 수의 노름이다. 또는 \(\mathfrak p\)-진 확대로 넘어가지 않고 말하면 다음과 같다. \emph{\(a\)가 \(k\)의 모든 소아이디얼 \(\mathfrak p\)에 대하여 노름잉여이고 몇 가지 부호조건도 만족하면, \(a\)는 \(Z\)의 한 수의 노름이다.} 이 마지막 형태가 바로 알려진 해석적 보조수단을 이용하여 유체론에서 증명되는 노름정리이다. 분할대수에 관한 일반정리는 이 순환적 특수 경우로부터 순전히 대수적·산술적 고찰로 증명할 수 있다.\footnote{R. Brauer, H. Hasse, E. Noether, \emph{Beweis eines Hauptsatzes in der Theorie der Algebren}, J. f. M. 167 (1932).} Hasse는 여기서 중요한 첫 귀결을 얻었다. \emph{대수적 수체 위의 모든 단순 정규대수는 순환대수이다.} 오래 추측되어 온 이 사실의 증명을 찾는 과정에서 위의 일반적 정식화가 생겨났다.

\paragraph{5.}
분할대수에 관한 정리의 두 번째 귀결은, 다시 순전히 대수적·산술적인 추론으로 얻어지는 처음의 \emph{주종정리}이다.\footnote{그 증명은 Math. Ann.에 실릴 예정이다.} 그 불변적 정식화는 교차곱을 정의하는 관계 (2)부터 (5)가 순전히 곱셈적이어서, \(K^\ast\)를 다음 조건만을 만족하는 아벨군 \(\mathfrak J\)로 바꾸어도 뜻을 유지한다는 사실에 바탕을 둔다. 즉 \(\mathfrak J\)의 자기동형군이 \(\Gg\)와 동형인 부분군 하나를 포함하면 된다. 이때 (1)은 군론에서 말하는 \glqq{}\(\mathfrak J\)에 의한 \(\Gg\)의 확장\grqq{}으로 바뀐다. \(\mathfrak J\)를 \(K\)의 모든 아이디얼로 이루어진 군으로 택하면 인자계도 아이디얼들의 계가 된다. \(\mathfrak J\)의 류분할은 인자계에 대한 류분할을 유도하는데, 일반적으로 이는 본래의 아이디얼류분할보다 더 세밀하다. \(u_S\)들의 곱셈이 절대 아이디얼류에 관하여 일의적이라는 요구는 변환량 \(c_S^{\,T}c_T/c_{ST}\), 곧 원소 인자계의 단위원류가 아이디얼 인자계의 단위원류에 속한다는 것만을 말하기 때문이다. 그러나 알려진 경우로 특수화해 보면 실제로는 이보다 조금 덜 세밀한 분할만으로 충분하다. 다음과 같이 정의한다. \emph{모든 유한 및 무한 분기소에서 분할되는 대수를 생성하는 원소 \(a_{S,T}\)들을 인자계의 단위원류에 넣는다.}

이렇게 생기는 \(\Gg\)의 확장을 \(\Gg^\ast\)로 쓰고, \(c\)의 절대 아이디얼류를 \((c)\)로 쓰자. 그러면 다음을 얻는다.

\emph{주종정리의 불변적 정식화: 위에서 정의한 류분할 아래에서 치환 \(v_S=u_S(c_S)\)가 \(\Gg^\ast\)의 자기동형사상을 만든다고 하자. 이러한 모든 \((c_S)\)가 주종을 이룬다. 그러면 그 자기동형사상은 내부자기동형사상이며, 하나의 아이디얼류 \((\mathfrak b)\)에 의해 생성된다.}

알려진 특수 경우들은 다음의 동치인, 조금 더 명시적인 정식화에서 따라 나온다. \emph{아이디얼류 \((c_S)\)로부터 만든 변환량 \((c_S^{\,T})(c_T)/(c_{ST})\)이 인자계의 단위원류에 속한다고 하자. 그러면 모든 \(S\in\Gg\)에 대하여 \((c_S)=(\mathfrak b)/(\mathfrak b^S)\)가 되도록 하는 아이디얼류 \((\mathfrak b)\)가 존재한다. 즉 \((c_S)\)들은 기호적 \((1-S)\)제곱이다.} 실제로 가정은 바로 자기동형사상 성질을 표현하고, 그 자기동형사상이 내부적이라는 사실은
\[
v_S=(\mathfrak b)^{-1}u_S(\mathfrak b)
   =u_S(\mathfrak b)^{1-S}
   =u_S(c_S)
\]
로 표현된다.

순환적인 경우로 특수화하면, 정규화를 고려하여 다음을 얻는다. \(N([c])\)가 인자계의 단위원류에 속하면 \((c)\)는 기호적 \((1-S)\)제곱, 곧 \((c)=(\mathfrak b)^{1-S}\)가 된다.

\paragraph{6.}
여기서 순환체와 이차형식에 관한 알려진 형태로 나아가기 위해, 먼저 이 정리가 모든 아이디얼류에 대하여 진술되었지만 통상적인 방식대로 분기소와 서로소인 아이디얼들로 제한해도 그 내용은 같음을 말해 둔다.

그러면 여기서 정의한 주종은 이차체의 경우 가우스의 주종으로 바뀐다. 아이디얼류에는 이차형식이, 류의 노름에는 그 형식들이 나타내는 수가 대응하기 때문이다. 단위원류가 \(K\)의 분기소에서 분할되는 대수들을 생성한다는 것은, 이 표현되는 수들이 분기소에서 이차잉여가 된다는 뜻이다. 따라서 대응하는 형식들은 주형식의 총특성(Totalcharakter)을 가지며 가우스의 주종을 이룬다. 기호적 \((1-S)\)제곱이 제곱연산으로 바뀐다는 사실은 잘 알려져 있다.

순환체의 경우에는 다음 형태가 된다. 주종은 그 노름이 모든 유한 및 무한 분기소에서 노름잉여가 되는 아이디얼류 전체로 이루어진다. 이는 잘 알려져 있듯이 \glqq{}도체에 대한 노름잉여\grqq{}와 동치이므로 통상적인 정리가 특수 경우로 얻어진다.

그 밖에도 임의의 갈루아체라는 일반적인 경우에 분기소들로만 이루어진 도체를 도입할 수 있다. 이때 각 분기소에서 인자계를 정규화하면 그 레이는 단위원류의 원소들만을 포함한다.

여기서 2절의 개관에서 말한 Artin 도체들과의 관계가 문제로 생긴다. Artin 도체들도 같은 소아이디얼들로 이루어져 있기 때문이다. 더 나아가 두 번째 초복소적 방법인 갈루아 가군론과의 관계도 문제가 된다. 이 두 방법이 어디까지 미칠지는 앞으로 밝혀져야 한다.

\endgroup
